\section{A Distinct Adjoint for Regional Thermal Protection}
\label{app:regional}

\subsection{Objective, adjoint, and the insertion score}
Fix a nonnegative weight $q$ supported in a protected region $S\Subset D$.
Require all admissible cavities to be disjoint from $S$ and set
$Q=\int_S q>0$. The physical source $f$ and the observation weight $q$
serve different roles. For the zero-Dirichlet Poisson state define
\begin{equation}
 -\Delta u_D=f\quad\hbox{in }D,\qquad
 J_q(D)=Q^{-1}\int_S q u_D.
 \label{eq:regionalobjective}
\end{equation}
Let $p_D$ solve $-\Delta p_D=q$ in $D$ with zero boundary data. This
unnormalized adjoint is convenient in the implementation; the actual
adjoint for $J_q$ is $p_D/Q$.

\begin{proposition}[Regional sensitivity and monotonicity]
Under the smoothness and separation assumptions of
\cref{thm:localtd}, the capacity-normalized derivative for a
zero-temperature cavity at $z\notin S$ is
\begin{equation}
 D_TJ_q(z)=-u_D(z)p_D(z)/Q.
 \label{eq:regionaltd}
\end{equation}
For finite nested domains obtained by removing cavities outside $S$,
$J_q$ is nonincreasing. This monotonicity does not establish a minimum
benefit, or certify a particular finite-radius prediction.
\end{proposition}
\begin{proof}
\textbf{Step 1: differentiate the objective in the state.}
For every $v\in H^1_0(D)$, the linear objective has derivative
$j'(u)[v]=Q^{-1}\int_D qv$. The adjoint weak equation is therefore
$\int_D\nabla v\cdot\nabla(p_D/Q)=Q^{-1}\int_D qv$.
The fixed support and cavity exclusion keep both $q$ and $Q$ unchanged
by every admissible insertion.

\textbf{Step 2: apply the cavity expansion with this adjoint.}
The state trace imposed on the inserted cavity is zero. Substitution
of $h_z=0$ and adjoint $p_D/Q$ into the established local coefficient
$-(u_D(z)-h_z)p(z)$ gives \cref{eq:regionaltd}. Multiplication by
the positive constant $Q$ does not change candidate rankings. This
specialization uses the classical coefficient; it is not a claim that
the underlying topological formula is new.

\textbf{Step 3: compare finite-domain states.}\par
Write $u_D^-=\max(-u_D,0)\in H^1_0(D)$. The weak equation gives
\[
 -\int_D|\nabla u_D^-|^2=\int_D f u_D^-\geq0.
\]
Both the gradient norm and the nonnegative right-hand side must vanish.
Poincar\'e's inequality and the zero trace imply $u_D^-=0$, hence
$u_D\geq0$. For $D'=D\setminus K$, the difference
$w=u_D-u_{D'}$ is weakly harmonic in $D'$, has zero outer trace, and
has nonnegative trace $u_D$ on the new cavity. Consequently
$w^-=\max(-w,0)$ has zero trace on all of $\partial D'$ and belongs
to $H^1_0(D')$. The harmonic weak equation gives
$-\int_{D'}|\nabla w^-|^2=0$. Poincar\'e's inequality again yields
$w^-=0$, or $w\geq0$. Hence
$J_q(D)-J_q(D')=Q^{-1}\int_S qw\geq0$. Applying this argument at
each insertion proves the finite nested-domain statement. It supplies
no positive quantitative lower bound beyond zero.
\end{proof}

\subsection{An exact counterexample to hottest-point placement}
\begin{proposition}[A hottest point need not maximize topological benefit]
On the unit ball in $\mathbb R^3$, take $f=1$ and
$q(x)=1+b x_1$ for $0<|b|<1$. For the unnormalized linear objective
$\int q u$, the hottest point is the origin, but the origin does not
locally maximize the insertion benefit $up$.
\end{proposition}
\begin{proof}
\textbf{Step 1: solve the state.} Writing $r=|x|$, the function
$u=(1-r^2)/6$ vanishes on the boundary and satisfies $-\Delta u=1$,
because $\Delta r^2=6$. Uniqueness of the weak Dirichlet solution
identifies this function as the state. It is strictly decreasing in
$r>0$, so its unique maximum is at the origin.

\textbf{Step 2: solve the adjoint.} The product rule gives
$\Delta(x_1r^2)=x_1\Delta r^2+2\nabla x_1\cdot\nabla r^2
=6x_1+4x_1=10x_1$. Consequently
\[
 p(x)=(1-r^2)\left(\frac16+\frac{b x_1}{10}\right)
\]
has zero trace and satisfies $-\Delta p=1+b x_1=q$.
Again uniqueness identifies the adjoint.

\textbf{Step 3: inspect the benefit at the hottest point.}
Direct multiplication gives
\[
 u(x)p(x)=(1-r^2)^2\left(\frac1{36}+\frac{b x_1}{60}\right),
 \qquad \partial_{x_1}(up)(0)=\frac b{60}\ne0.
\]
Moving a sufficiently small distance in the direction
$\operatorname{sign}(b)e_1$ therefore increases $up$, while it
strictly decreases $u$. Thus hottest-point ranking and topological
ranking are distinct even for the same self-adjoint Poisson operator.
This analytic example establishes a mechanism, not the size of the
benefit in the housing experiment.
\end{proof}

\subsection{Paired random walks and what cross-fitting removes}
At a candidate, one walk accumulates both the physical-source
contribution $U_i$ and the observation-source contribution $P_i$.
The spherical exits and Green-distributed source samples are shared
within a pair. Different complete pairs use separate random streams.
For $N\geq2$ the implementation evaluates
\begin{equation}
 \widehat s=\frac{(\sum_iU_i)(\sum_iP_i)-\sum_iU_iP_i}{N(N-1)}.
 \label{eq:regionalcross}
\end{equation}
\begin{proposition}[Expectation of the implemented product]
For independent identically distributed pairs with finite second
moments, write $\mu_U=\E U_i$, $\mu_P=\E P_i$. Then
$\E\widehat s=\mu_U\mu_P$, whereas
$\E[N^{-1}\sum_i U_iP_i]=\mu_U\mu_P+\Cov(U_i,P_i)$.
\end{proposition}
\begin{proof}
Expanding the product of sums and subtracting its diagonal leaves
$\sum_{i\ne j}U_iP_j$. There are exactly $N(N-1)$ ordered distinct
pairs. For each, independence gives
$\E[U_iP_j]=\E U_i\E P_j=\mu_U\mu_P$. Divide the sum by its
number of terms to obtain the first assertion. For the diagonal,
the defining identity for covariance is
$\E[U_iP_i]=\mu_U\mu_P+\Cov(U_i,P_i)$; averaging preserves it.
\end{proof}
Cross-fitting removes product-correlation bias. It does not remove
WoS termination error. In the ideal contained-ball model with
nonnegative sources, dropping the terminal values makes each
expectation a lower approximation of its state. The measured code
uses the existing approximate distance and finite-precision sampling
implementation and does not assert a rigorous bias magnitude.
The specialized compliance certificate cannot be substituted into
\cref{eq:regionalobjective}: its square-energy identity is specific
to compliance. These new designs use fresh ranking at a fixed count.
Radius-independent variance does not imply independence from the
observation region: a very small support of $q$ can itself make source
sampling difficult. The regional study fixes that support before varying
the estimator budget. Its source-sampling difficulty is distinct from
encountering a shrinking trial cavity.

\subsection{Sensor normalization and protected geometry}
The two observation regions are radius-$a=0.13$ balls centered at
$(\pm0.55,\pm0.18,-0.03)$, with the signs chosen together. Define
$q(x)=(1-|x-c|^2/a^2)_+^2$. Radial integration gives
\begin{align*}
 Q&=4\pi a^3\int_0^1 t^2(1-t^2)^2\,dt\\
  &=4\pi a^3\left(\frac13-\frac25+\frac17\right)
    =\frac{32\pi a^3}{105}.
\end{align*}
The sensor-ball volume is $4\pi a^3/3$, so its average weight is
$8/35$. Uniform sensor-ball probes therefore estimate $J_q$ as the
sample mean of $qu$ divided by $8/35$. Candidate centers obey
$|z-c|\geq a+r+0.02$; by the triangle inequality every point of the
radius-$r$ cavity is at least $0.02$ away from the protected ball.
Thus reducing the objective cannot be accomplished by deleting
the monitored region.

\subsection{Replicated designs and numerical reference checks}
\begin{table*}[t]
\centering\small
\caption{Regional design comparison. Reductions are means $\pm$ one standard error across five designs per sensor. Both wall-time columns include candidate preparation and the complete selection driver, averaged across sensor cases. Retained-mesh timings use one fresh server per design; dashes denote configurations not rerun with that backend. The extra multigrid resolutions are a supplementary sweep.}
\label{tab:regional}
\begin{tabular}{lrrrr}
\toprule
Method & Left reduction (\%) & Right reduction (\%) & Reload mesh (s) & Retain mesh (s) \\
\midrule
Random & $10.55\pm1.54$ & $10.23\pm3.14$ & 0.44 & 0.39 \\
WoS hottest & $36.47\pm2.22$ & $29.26\pm2.24$ & 3.63 & 1.19 \\
MG hottest & $38.15\pm1.03$ & $29.99\pm1.77$ & 6.07 & 4.88 \\
TWoS sparse & $50.87\pm1.46$ & $50.09\pm1.34$ & 3.66 & 1.22 \\
TWoS larger budget & $52.13\pm1.08$ & $51.84\pm0.52$ & 5.04 & 2.62 \\
MG TD (57) & $53.50\pm0.62$ & $52.53\pm0.84$ & 2.51 & 1.41 \\
MG TD (85) & $53.57\pm0.66$ & $52.45\pm0.82$ & 4.54 & -- \\
MG TD (113) & $53.30\pm0.57$ & $52.15\pm0.88$ & 8.38 & 7.13 \\
MG TD (169) & $53.57\pm0.66$ & $52.15\pm0.88$ & 40.00 & -- \\
\bottomrule
\end{tabular}
\end{table*}

The benchmark retains all final geometries, per-round candidate fields,
solver residuals, work counts, and timings. Five seeds change the shared
candidate set and stochastic streams. The two target regions are evaluated
separately on the same housing. Each final objective uses four independent
scrambled Sobol blocks of 128 uniform sensor-ball probes, with 256 fresh
WoS paths per probe. Evaluation streams are independent of selection and
shared between methods for paired comparisons. The table reports variation
across complete designs, not a confidence guarantee for arbitrary new inputs.
The separate reference solver evaluates the state at the same probes.
Across first-seed designs, the fourfold shell refinement changes measured reduction by at most $0.087$ percentage points. Refining the evaluation grid from $169\times109\times85$ to $225\times145\times113$ changes it by at most $0.255$ points. Across all primary and supplementary designs, WoS and the $169\times109\times85$ reference differ by at most $1.864$ points. These sensitivity checks do not bound the remaining discretization or sampling bias.


The primary design protocol and sampling budgets were fixed before the
first completed seed. That seed was inspected before the other four ran.
The supplementary grid resolutions were specified after that pilot and
before their measurements. All are retained. No improvement in the
cooling model or guarantee of accessible cooling follows from changing
the observation objective. The theoretical counterexample above and
the observed improvement over hottest-point placement show why temperature
alone need not characterize objective sensitivity. The adjoint supplies the
missing observation dependence in the topological formula; deterministic
topological derivatives exploit the same information.
